\chapter[What Should Liberal Education Do?]{What Should a Liberal Education Finally Teach?}
\section{An Old Book Filled with Strangers' Handwriting}
First, pick up an old book.

Not a new one, but something found in a secondhand shop or a school library corner. Its cover frays, half its spine lettering has worn away, and opening it releases the smell of damp paper dried again. What matters is inside: writing fills the margins.

Not one person's writing. Beside Chapter 3, a blue-black fountain pen says, ``The author loses confidence here. The next page proves it.'' Narrow, forceful strokes suggest a young hand. In Chapter 7, faint pencil in another hand answers, ``Disagree with the person above. He's simply honest.'' Farther on, three heavy underlines and an exclamation mark sit below a page whose top says, ``Boring. Skip.''

You do not know these people. They probably never met, perhaps lived decades apart. Yet their marginal discussion has no moderator: judgments, rebuttals, and marks await later readers' encounters.

The book now has a strange property. It is no longer solely its author's work but strangers conversing through time. You, the next reader, must decide whether to read the notes. When they oppose your views, do you cross them out, argue back, or first ask why someone thought that?

This book in miniature contains the humanities course. Liberal education, your broad reading and thinking across subjects, should finally teach not accumulated titles and dates but what sort of person faces strangers' notes. We will unpack that question.

\section{The Last Reading Lesson}
Enter a scene with me.

Time: June's final afternoon lesson. Place: a ninth-grade classroom, where a creaking fan cannot shift the heat. Cicadas sound outside; half an unerased formula remains from math.

This is the year's final reading class. Ms. Lin brings a stack of poetry collections and places them face down. The class has read this book for a whole semester, learning more than a dozen poems, many by heart. Some quoted them in essays; others copied lines onto desk corners.

Last week, news broke. Recently published private letters by the long-dead poet contained cruelty: belittling women, mocking poor people, repaying supportive friends with ingratitude. Reports spread everywhere and online arguments erupted.

Instead of lecturing, Ms. Lin writes, ``Should this poetry collection remain on next year's reading list?'' She invites discussion.

The room explodes.

The Chinese class representative rises first from the third row, voice tight: ``Remove it. The letters show his real thoughts. We spent a semester moved by love poems while he secretly mocked people like us. That's fraud, not art.''

Her deskmate disagrees: ``Poems and people are different. Good writing doesn't depend on its author's character. By that standard, museums would be half empty.''

A usually quiet boy mutters from the back: ``But knowing those letters changes how the poem tastes. I can't help it. It just feels wrong.''

Someone in a corner raises a hand: ``Shouldn't we check whether all the letters are verified? This wouldn't be the first publication of private letters stripped of context.''

Ms. Lin remains silent, recording names and reasons. By the bell, the board is full, and a vote would probably split three to three to three: remove, retain, uncertain.

Remember this dispute, especially the boy whose altered feeling is a real difficulty rather than an argument. Any complete answer must make room for it.

\section{One Problem Is Really Six}
Do not vote yet. The single question of removal entangles at least six, precisely the capacities liberal education ultimately addresses.

\textbf{First: can you understand strangers' reasons without immediately agreeing?} Everyone thinks themselves right and others wrong, yet understanding differs from agreement. You can restate a deskmate's case to their satisfaction and still reject it. Most people lack the first half: preparing rebuttals by the third sentence, they answer an imagined opponent rather than the real one.

\textbf{Second: can you respect tradition without surrendering criticism?} Earlier generations read, loved, and debated this collection, lending traditional endorsement. Should that count, and how much? Overturning tradition and kneeling before it are easy. Standing upright without overturning the table is harder.

\textbf{Third: can you love a work without worshiping its author?} Good writing quietly becomes good person, then everything they say is right, until revelation shatters the structure. Where is the boundary between love and deification?

\textbf{Fourth: can you admit uncertainty without concluding we know nothing?} Verification and missing context matter, but doubt does not mean everything is false. Is there space between certainty and abandoning judgment?

\textbf{Fifth: can you judge and revise when evidence changes?} If archivists prove key letters forged tomorrow, is today's remove vote changing wisely or inconsistently? How do we treat changed minds?

\textbf{Sixth, deepest: what is this learning for?} Reciting poets' names at dinner to seem learned, or living more responsibly with disagreeing people in classrooms, cities, and networks?

None has a ready-made answer, but all are trainable capacities. Take them in turn.

\section{State Every Side at Its Strongest}
Before judging three positions, \textbf{make each argument fully enough for its supporters to recognize it}. This is steelmanning: build the strongest version before rebutting. Its opposite, the straw man, distorts a position into something easily demolished. Satisfying though that victory feels, it defeats air.

\textbf{First: remove it. Character and work cannot separate.}

The representative's fullest case has four layers. Work emerges from a person's mind, not the sky; if part is rotten, which lines carry the smell? Reading entrusts emotion. Discovering the author mocked readers like you is genuine betrayal, justifying withdrawn trust. A reading list is official recommendation, not neutral storage. Retaining a corrupt person's work endorses them and tells future students talent excuses conduct. Finally, good poems are abundant and shelf space limited. What great loss comes from giving the place to an overlooked, untainted author?

This strong case has an old intuition behind it: Chinese criticism says calligraphy and prose resemble their makers. For millennia, people have linked what appears on paper to the person behind it.

\textbf{Second: retain it. Once written, the work belongs to the world.}

The deskmate's fullest case also has four layers. A poem's language, rhythm, and imagery lie visibly open to examination, needing no character warranty. Authors cannot control completed works: a thousand readers find a thousand meanings, unmoved by declarations that all are wrong. If meanings escape the author, why subject them to collective moral punishment? Character screening also creates disaster: who sets standards, how far back, and do private letters count as works? Few historical writers may survive. Finally, tightly binding work and person nourishes both worship and collapse: elevating authors through work, then burning work through authors' flaws repeats the same refusal to judge them separately.

\textbf{Third: read with the full account. Neither remove nor pretend.}

Make the unsettled boy's case. He lacks theory but faces a genuine jolt when once-tender words recall cruel letters. Neither side should dismiss that. Keep the collection, but bring the letters into class too: who wrote, what they did, why some praise the poem and others condemn the person. Let students decide, building neither shrine nor effigy. The cost is trouble: more teacher preparation and students holding beautiful poetry and bad character together without deleting either. Perhaps bearing that trouble is education's purpose.

Notice an easily misread counterexample. If the representative accurately restates the retention case, someone may accuse her of helping opponents. \textbf{Strengthening another's argument is often mistaken for changing sides}. Her vote may remain removal. Understanding and agreement differ; confusing them begins discussion's decline into allegiance. Liberal education's first practice stores them separately.

This is not one classroom's problem. Ezra Pound was arrested in the 1940s for broadcasts supporting Italian fascism, yet received a leading American poetry prize while detained. Judges broadly defended judging poems rather than people, and dispute persists. Nazi Party member Martin Heidegger remains widely required in philosophy, equally controversially. A. K. Ramanujan's essay surveying hundreds of South Asian Ramayana versions entered Delhi University's curriculum, then was removed after protests in 2011. Opponents saw injury to sacred tradition; supporters saw recognition of diversity as respect. Person and work, faith and scholarship, respect and criticism collide across civilizations under different names.

\section[Add Confidence Levels to Judgments]{Thinking Bridge: Add Confidence Levels to Judgments}
The student asking about verification already holds a portable tool: a \textbf{confidence scale}.

Confidence is how firmly you hold a judgment. Instead of only yes and no, install at least four levels:

\begin{itemize}
\item \textbf{I am certain}: direct evidence or repeatedly tested consensus. The collection was on last semester's list; you read it yourself.
\item \textbf{I think it highly likely}: reliable sources support it, without personal verification. Several serious outlets report the letters' publication.
\item \textbf{I suspect}: some indications, thin evidence. The letters reflect lifelong beliefs, perhaps, or only an angry difficult period.
\item \textbf{I do not know}: insufficient grounds even to guess. The poet's thoughts while composing that love poem may have been unclear even to him.
\end{itemize}
The crucial next step is attaching \textbf{what evidence would change my level} to each judgment.

\begin{itemize}
\item Certain the letters were published: a publisher or archive demonstrating forgery could overturn it.
\item Suspect the letters reflect genuine beliefs: a complete collection rather than extracts might reveal different context.
\item Do not know the poet's thoughts: evidence to upgrade may never arrive. Leaving it unknown is no disgrace.
\end{itemize}
The scale blocks two errors. \textbf{All or nothing} demands complete certainty or believes nothing. \textbf{Confidence drift} turns yesterday's guess into common knowledge after three forwards. Most online heat comes from first-level tones backed by third-level evidence.

Apply it to the classroom board. Much disagreement occurs across levels. Remove and retain involve legitimately debatable evaluation, while authenticity is factual and presently likely but unverified. Before it settles, conclusions hang provisionally. Asking the level before arguing saves much breath.

This tool guarantees not correctness but \textbf{honesty} about knowledge and ignorance. The confidence needed for action is a further question.

\section{The Sentence Mencius Throws at Us}
Read a primary source. Someone spoke firmly about critical respect for tradition 2,300 years ago.

In Mencius's ``Jin Xin II'':

\begin{quote}
Better no Book of Documents than believing all of it.
\end{quote}
The Book of Documents was an authoritative record of ancient rulers, comparable to bold, starred textbook knowledge today. Mencius says believing everything makes the book worse than none. He accepts only twenty or thirty percent of ``Successful Completion of the War,'' rejecting its blood-floating-pestles description of King Wu's campaign against Zhou as unreasonable. His benevolent ruler should not kill rivers of blood.

Notice both sides. Mencius devoted his life to defending ancient kings' tradition, not rejecting it. Yet he \textbf{weighs classics with his own reason} rather than placing them beyond touch. He accepts one passage and only thirty percent of another when reasonableness fails. Tradition counts, but judgment stays with him.

Pair this with the Analects' ``Wei Zheng'':

\begin{quote}
Learning without thought leaves one bewildered; thought without learning leaves one in danger.
\end{quote}
Learning without reflection yields confusion; unsupported speculation lacks foundations. Like pincers, the two lines constrain us. Criticism without respect attacks imagined tradition one never studied; respect without criticism lets others' horses run through one's mind.

Admit Mencius's difficulty too. His thirty-percent standard is his own benevolence, itself inherited from tradition. Does using tradition to judge tradition merely retain what pleases and dismiss what does not? Later readers criticized this bold interpretive projection. Critics are shaped too. Nobody stands wholly outside tradition; we can know and disclose our positions.

Traditions also contain many voices. Medieval Islam saw al-Ghazali's Incoherence of the Philosophers criticize overreaching Greek reasoning, then Ibn Rushd's Incoherence of the Incoherence defend rational inquiry point by point decades later. Both inhabited one faith, took each other seriously, and yielded little. Tradition is a millennia-long debate, not a neatly sealed box. Joining means taking a microphone, not accepting a parcel.

\section{One Poet, Three Readings}
Return with more tools. Separate what happened, publication and reporting, from explanations of the shock, participants' accounts, and judgments about response. Here compare positions within the latter two levels, each distinct and limited.

\textbf{Reading one: the autonomous work.} Completed art becomes independent, judged by its own language and artistry. Mid-twentieth-century Anglo-American New Criticism's intentional fallacy treats appeals to authorial intention and biography as misdirection. Letters therefore cannot determine poetic quality, making removal unwarranted. This protects works from shifting moral winds but pretends reading is sterile. It can tell the unsettled boy not to feel that way, not cure his feeling.

\textbf{Reading two: moral connection.} Art extends personality, consumption relates readers to authors, and official recommendation sustains the discredited through institutional authority. Removal expresses a school position rather than censorship. This takes endorsement and injured feelings seriously, but screening lacks a stopping line: who, by which standard, how deeply? Each question grows harder, and historical abuse exceeds beneficial use.

\textbf{Reading three: informed reading.} Retain the work without hiding its maker. Teach controversy and let informed students hold good work and bad author together. Crisis becomes education, training decision-makers rather than deciding for them. It demands much from teachers and students and attracts both sides' complaints: evasive compromise to removal supporters, unnecessary complication to retention supporters.

We announce no winner. Their deepest differences concern who owns meaning, what recommendation's authority signifies, and whether education filters the world for young people or accompanies their practice in facing it.

\section[Fallibilism: Firm but Revisable]{Deeper Challenge: Fallibilism, Firm but Revisable}
Our strongest theory addresses admitting uncertainty without nihilism and judging while remaining revisable.

\textbf{Fallibilism}, systematically developed in the twentieth century by Karl Popper, begins with two apparently conflicting statements. Any human knowledge may be wrong, with no ultimate guarantor; authorities, classics, and consensus have been overturned. Yet knowledge exists and grows: today's greater understanding than Mencius's era is no illusion. Popper studied science, but the attitude reaches every judgment.

How coexist? \textbf{Knowledge matters not because it is proven eternally correct, but because it withstands and welcomes tests}. A claim specifying how the world would look if it were wrong deserves serious attention. One immune to every outcome or evidence falls outside knowledge. In the preface to Conjectures and Refutations, Popper broadly expresses an attitude: I may be wrong, you may be right, and together we may approach truth.

Compare its two opponents; these distinctions are central.

\begin{itemize}
\item \textbf{Dogmatism}: I am certain; leave me alone. Certainty ends inquiry.
\item \textbf{Radical skepticism}: anything may be wrong, so believe nothing and judge nobody. Possible error excuses surrender.
\item \textbf{Fallibilism}: possible error makes reasons, evidence, and objections more necessary. My present judgment deserves care precisely because it is my best current version, not an indifferent whatever.
\end{itemize}
Skepticism and fallibilism share possible error but move toward resignation or diligence. Two people knowing they will die likewise conclude either life is pointless or today deserves care. The premise matches; attitude intervenes between opposite conclusions.

Ancient roots remain. In Plato's Apology, Socrates recounts an oracle naming him wisest. Doubting it, he visits reputedly wise politicians, poets, and craftsmen and finds misplaced confidence outside their knowledge. His small advantage is not imagining he knows what he does not. He does not conclude nobody knows anything; awareness of ignorance fuels lifelong questioning. It is an engine, not a brake.

Answer a classroom challenge. Someone believed last year and disbelieves now; we call it flip-flopping, humiliation, or opportunism. An easily misread counterexample appears: \textbf{cultures praising stable positions punish evidence-based revision}. Fallibilism reclassifies changed minds by reasons, not conclusions. New evidence or collapsed arguments mean a functioning system, not betrayal. Changing with pressure or fashion is the real turncoat behavior. Ask evidence first, prevailing winds second.

Admit the limits. Fact judgments such as authenticity admit tests more readily than values such as whether a poem should move someone. People are not argument machines. Trust in friends and community commitments sometimes require standing without sufficient evidence; perpetual readiness to let go can become evasion. Fallibilism is an excellent map, not the territory. Love and loyalty blur at its scale.

\section{Worship, Collapse, and Information Cocoons}
Every question here plays daily on phones, amplified tenfold by algorithms.

\textbf{Worship and collapse are opposite ends of one production line.} Fandom turns good art into good character, then universal correctness, with questioners treated as enemies. The more perfect worship becomes, the more devastating scandal feels, because believers have only god and garbage, with no slot for a complex, flawed maker of good work. The classroom repeats across the internet, but without Ms. Lin. Half the pain is interest on refusing to separate person and work.

\textbf{Algorithms select your strangers.} Liberal education asks you to hear unfamiliar reasons; recommendations instead supply what you enjoy. Opposing views arrive in their foolishest forms because you enjoy their embarrassment. You think you learn opponents daily while defeating straw men and gaining confidence. The technically simple but difficult remedy is actively seeking their best argument, not their most viral one, and steelmanning before rebuttal. Algorithms will not do this for you; education still must.

\textbf{Uncertainty suffers two distortions online.} One voice declares certainty about everything, treating feelings as verified facts. Another calls everything fake because experts disagree, sliding from possible error into giving up. Opposites share unwillingness to do the difficult middle work: mark confidence, verify, suspend provisionally. Online fallibilism is the scholarly name for thinking three seconds before sharing.

\textbf{Changing minds carries a high online price.} Three-year-old posts await screenshots. Public positions incur humiliation penalties for revision, encouraging hidden thoughts or stubbornness until positions freeze and private thought stops. Healthy discussion should award medals for evidence-based changes, not screenshots. Participate now: when someone carefully explains a former mistake, inspect whether evidence persuaded them or fashion pushed them before mocking.

The final point is that \textbf{humanities learning is not for appearing learned}. Knowing Pound's case, memorizing Mencius, or defining fallibilism has no value by itself: these are signs toward actions. Hear an uncomfortable person out; ask confidence before forwarding; love a poem without building a temple; respect tradition while retaining judgment; form a view and watch what evidence could change it. Without these acts, reading is warehouse inventory. Their shared destination is responsible coexistence with people who will disagree throughout your life in classrooms, cities, and networks. These practices are how we handle that inevitability.

\section[How Would You Change Your Mind?]{For You: How Would You Want to Change Your Mind?}
Return to the classroom.

The bell has rung, writing remains, and books lie face down. Ms. Lin assigns one sentence: write a paragraph on retaining the collection next year; she will judge reasons, not conclusions.

Now your turn, with an enlarged assignment connecting everything from Chapter 6 onward:

\textbf{Imagine someone finds today's paragraph ten years later and discovers you no longer agree. What would you want to have changed your mind?}

New evidence confirming or exposing forged letters? A stronger argument revealing a gap? Experience making the boy's altered feeling understandable? Or everyone around you changing while standing alone became too tiring?

The first three probably deserve respect; the fourth probably does not. Yet how do you distinguish evidence's persuasion from a crowd wearing you down while inside it? Popper offers direction, not a lie detector. Only you know why you changed, and even you may deceive yourself.

For a question without a standard answer, a similarly open suggestion: occasionally record important judgments and reasons from today onward. Like strangers in old margins, leave a note for a reader ten years later, yourself. The old annotated book's final page is blank.

It is reserved for you.

\backmatter
